% 王江讲义树图候选样式；由正式入口加载 tikz、calc、arrows.meta 后接入。
% 单期接口兼容既有 \WJEventTree；接入时替换旧定义，勿重复定义。
\makeatletter
\newif\ifWJ@tree@stem
\pgfkeys{
  /WJ tree/.is family,
  /WJ tree/left/.initial={},
  /WJ tree/stem/.code=\WJ@tree@stemtrue,
  /WJ tree/row sep/.initial=1.7em,
}
\NewDocumentCommand{\WJEventTree}{O{} m}{%
  \begingroup
  \WJ@tree@stemfalse
  \pgfqkeys{/WJ tree}{#1}%
  \begin{tikzpicture}[x=1em,y=1em,baseline=(current bounding box.center)]
    \def\WJ@n{0}%
    \foreach \WJ@lab [count=\WJ@i] in {#2} {\xdef\WJ@n{\WJ@i}}%
    \pgfmathsetlengthmacro{\WJ@dy}{\pgfkeysvalueof{/WJ tree/row sep}}%
    \pgfmathsetlengthmacro{\WJ@ymax}{\WJ@dy*(\WJ@n-1)/2}%
    \draw[line width=.4pt,line cap=butt] (0,\WJ@ymax) -- (0,-\WJ@ymax);
    \foreach \WJ@lab [count=\WJ@i] in {#2} {%
      \pgfmathsetlengthmacro{\WJ@y}{\WJ@ymax-\WJ@dy*(\WJ@i-1)}%
      \draw[line width=.4pt,line cap=butt] (0,\WJ@y) -- (.7,\WJ@y);
      % mid west 将数学轴置于支线高度；不再使用 base west 和手工 yshift。
      \node[anchor=mid west,inner xsep=.28em,inner ysep=0pt]
        at (.7,\WJ@y) {$\WJ@lab$};
    }%
    \pgfkeysgetvalue{/WJ tree/left}\WJ@left
    \ifx\WJ@left\empty\else
      \WJ@tree@stemtrue
      \node[anchor=mid east,inner xsep=.22em,inner ysep=0pt]
        at (-.85,0) {\WJ@left};
    \fi
    \ifWJ@tree@stem
      \draw[line width=.4pt,line cap=butt] (-.85,0) -- (0,0);
    \fi
  \end{tikzpicture}%
  \endgroup
}
\makeatother

% 多期二叉树：节点为先、连线为后；不预设重合或路径展开结构。
% 参考 431《金融衍生品》具名节点和按时点分列的布局，本项目独立维护。
\tikzset{
  WJ binomial/.style={font=\small,baseline=(current bounding box.center)},
  WJ binomial compact/.style={WJ binomial,x=2.3cm,y=.78cm},
  WJ binomial node/.style={align=center,inner xsep=2pt,inner ysep=1pt,
    minimum width=7mm,minimum height=5mm,outer sep=1pt,
    execute at begin node={\setlength{\baselineskip}{10.5pt}}},
  WJ binomial edge/.style={-{Stealth[length=1.5mm,width=1mm]},line width=.45pt},
  WJ binomial time/.style={font=\small,anchor=north,inner sep=1pt},
}
